%قبل از شروع مسئله ترجیحاً چیزی ننویسید. اما اگر خواستید بنویسید، در قالب لیستی (فقرات و شمارش) بنویسید.
%برای اضافه کردن نمره، بعدی از شروع‌مسئله، [x] را اضافه کنید و به جای x نمره مورد نظر را قرار دهید. برای زیرمسئله‌ها هم به همین صورت است.
%\شروع‌مسئله[10]
\شروع‌مسئله [20]
%\متن‌سیاه{سوال عملی اول:}\\
%\متن‌سیاه{دراپ‌تست}\\

% صورت اصلی مسئله را در این‌جا بنویسید.
جایگشتی به طول $n$ داریم، هدف ما مرتب‌سازی این جایگشت است. ۲ نوع عملیات داریم:


۱. ۲ عدد متوالی در جایگشت را در نظر بگیریم و جایشان را عوض کنیم.

۲. ۳ عدد متوالی از جایگشت را در نظر بگیریم و ترتیبشان را برعکس کنیم.

می‌خواهیم طوری این جایگشت را مرتب کنیم که تعداد دفعاتی که از عملیات نوع ۱ استفاده می‌کنیم کمینه شود (تعداد دفعات استفاده از عملیات نوع ۲ اهمیتی ندارد).
در تنها خط خروجی حداقل تعدادی که می‌توانیم از عملیات نوع ۱ استفاده کنیم و جایگشت را مرتب کنیم چاپ کنید.

\متن‌سیاه{ورودی:}\\
%در خط اول ورودی، ابتدا k و سپس m داده می‌شود.\\
در خط اول ورودی عدد $n$
و
در خط دوم یک جایگشت از اعداد $1$ تا $n$ به شما داده می‌شود.

\متن‌سیاه{خروجی:}\\
%در تنها سطر خروجی، تعداد آیفون‌های مورد نیاز می‌آید. اگر در هیچ صورتی چنین کاری ممکن نبود ۱- چاپ می‌شود.\\
در تنها خط خروجی حداقل تعداد استفاده از عملیات نوع ۱ را چاپ کنید.

\متن‌سیاه{محدودیت‌ها:}\\
$1\leq n\leq 10^5$

%$m\leq10^9$

%\textbf{\متن‌سیاه{ورودی و خروجی نمونه}\\
\متن‌سیاه{ورودی و خروجی نمونه:}\\


%% تست اول
\begin{latin}
\begin{tabularx}{\textwidth}{|X | X|}
\hline
\textbf{stdin} &
\textbf{stdout}\\
\hline
5

1 2 5 4 3&
0\\
\hline

\end{tabularx}
\end{latin}

%% تست دوم

\begin{latin}
\begin{tabularx}{\textwidth}{|X | X|}
\hline
\textbf{stdin} &
\textbf{stdout}\\
\hline
5

1 3 5 4 2&
1\\
\hline

\end{tabularx}
\end{latin}

%اگر مسئله چند قسمت داشت به این صورت عمل کنید:
%\شروع{زیرمسائل}
%\زیرمسئله[0] % مسئلۀ ۱آ
%\زیرمسئله % مسئلۀ ۱ب
%\زیرمسئله % مسئلۀ ۱پ
%\زیرمسئله % مسئلۀ ۱ت
%\پایان{زیرمسائل}

\solution{
عملیات نوع ۲ مانند این است که ۲ عدد که در جایگشت فاصله‌شان ۲ است را جابجا کنیم. با این عملیات می‌توانیم به خانه‌هایی که در جایگاه فرد هستند و خانه‌هایی که در جایگشت زوج هستند هر ترتیبی بین خودشان بدهیم.

حال فقط کافی است اعداد فرد را به جایگاه‌های فرد و اعداد زوج را به جایگاه‌های زوج ببریم، برای این کار هر مرحله ۱ عدد زوج که در جایگاه فرد است را به خانه اول می‌بریم، و یک عدد فرد که در جایگاه زوج است را به جایگاه دوم می‌بریم و جای این ۲ عدد را با هم عوض می‌کنیم.

در این‌صورت تعداد عملیات‌های نوع ۱ که استفاده می‌کنیم برابر خواهد بود با تعداد عدد‌های فرد در خانه‌های زوج.

به راحتی می‌توانید ببینید که این تعداد لازم است و با کمتر از این تعداد نمی‌توانیم.


}




%%%%%%%%%%%%%%%%%%%%%%%%%%%%%%%%%%%%%%%%%%%%%%%%%%%%%
% برخی دستورات مفید برای تمرین‌ها. %
%%%%%%%%%%%%%%%%%%%%%%%%%%%%%%%%%%%%%%%%%%%%%%%%%%%%%

% فرمول‌های نمونه:
%$f(x) = \Theta(x)$
%$T(x, y) \leq \log(x) + 2^y + \binom{2n}{n}$
% {\tt code\_function}


% لیست موارد بدون شماره‌گذاری:
%\شروع{فقرات}
% \فقره مورد اول 
% \شروع{فقرات}
% \فقره مورد اول داخلی 
% \شروع{فقرات}
% \فقره اولین مورد خیلی داخلی 
% \فقره دومین مورد خیلی داخلی 
% \پایان{فقرات}
% \فقره مورد دوم داخلی 
% \پایان{فقرات}
% \فقره مورد دوم 
%\پایان{فقرات}

% لیست موارد شماره‌گذاری شده:
%\شروع{شمارش}
% \فقره یک مورد 
% \فقره یک مورد دیگر 
% \فقره باز هم یک مورد دیگر
%\پایان{شمارش}

% برای نوشتن یک فرمول چند خطی که یک شماره بخوره. برای حالت‌های دیگر به سایت زیر مراجعه کنید:
% https://www.overleaf.com/learn/latex/Aligning_equations_with_amsmath
%\begin{align} 
% \begin{split}
% (x+y)^3 &= (x+y)^2(x+y) \\
% &= (x^2+2xy+y^2)(x+y) \\
% &= (x^3+2x^2y+xy^2) + (x^2y+2xy^2+y^3) \\
% &= x^3+3x^2y+3xy^2+y^3
% \end{split} 
%\end{align}

% نوشتن ماتریس و آرایه دوبعدی
%\begin{align}
% A = 
% \begin{bmatrix}
% A_{11} & A_{21} \\
% A_{21} & A_{22}
% \end{bmatrix}
%\end{align}

% اضافه کردن تصویر با شرح. برای اضافه کردن تصویر بدون شرح، فقط دستور درج‌تصویر را استفاده کنید.
%\شروع{شکل}[h!]
%\تنظیم‌ازوسط
% \درج‌تصویر[پهنا=0.5\پهنای‌متن]{img.jpg}
%\شرح{شرح تصویر}
%\برچسب{تصویر:۱}
%\پایان{شکل}
%\پایان{زیرمسائل}

%کد زیر را ببینید.
%%\lstinputlisting{source_filename.py}
%\begin{latin}
%\begin{lstlisting}
%# Import the modules
%import sys
%import random
%
%ans = True
%
%while ans:
%question = raw_input("Ask the magic 8 ball a question: (press enter to quit) ")
%
%answers = random.randint(1,8)
%
%if question == "":
%sys.exit()
%
%elif answers == 1:
%print "It is certain"
%
%elif answers == 2:
%print "Outlook good"
%
%elif answers == 3:
%print "You may rely on it"
%
%elif answers == 4:
%print "Ask again later"
%
%elif answers == 5:
%print "Concentrate and ask again"
%
%elif answers == 6:
%print "Reply hazy, try again"
%
%elif answers == 7:
%print "My reply is no"
%
%elif answers == 8:
%print "My sources say no"
%\end{lstlisting}
%\end{latin}